% ==================================================================
%  LabelPropagation.tex  --  Section 3.2.3 "Label Propagation"
%  Loaded from main.tex.
%  Implementation: Label_Propagation.py, fhe_kd_dense_propagation
%  (_slot_group_size, _pack_candidates, _tree_max_packed).  Max calls per
%  round: 2 * sum_j ceil(log2(|G_j|+1)).
% ==================================================================

\subsubsection{Label Propagation}
\label{sec:lp}

Label Propagation is the final step on the server. It turns the two
masks into cluster labels. The original DBSCAN~\cite{ester1996density}
does this with a traversal: it picks an unvisited core point, follows the
chains of neighboring core points, and gives every point it reaches the
same label. A traversal is out of reach here for two reasons. The server
cannot tell which point comes next, since that depends on values it holds
only in encrypted form. And even if it could, the work would differ from
point to point---one point expanded into many neighbors, another into
none---so the time and the number of operations spent around a point
would let the server infer the distribution of the data.

Distributed variants of DBSCAN replace the traversal by propagation.
NG-DBSCAN~\cite{lulli2016ng}, for instance, tags each cluster with
the label of its highest-ranked node, ranking nodes by their degree. We
follow the same max-based propagation and rank the core points by their
slot. Every core point starts with a label of its own, labels spread
between neighboring core points, and wherever two labels meet the larger
one is kept. Over the rounds, the core points of a cluster come to carry
the same label, the largest they started with.

\paragraph{\emph{Labels and the Maximum}}
The initial label of slot~$i$ is $u_i = (i+1)/n$. Core points keep it
and all other points start with $0$:
\begin{equation}
  \label{eq:init}
  \boldsymbol{\ell} = \mathsf{Core\_mask} \odot \mathbf{u},
  \qquad u_i = \frac{i+1}{n} ,
\end{equation}
where $\mathbf{u}$ is a plaintext vector. Dividing by $n$ puts every
label in $[0,1]$ and every difference of two labels in $[-1,1]$. A
difference can reach the ends of that interval, for instance when the
label $1$ meets a candidate that a mask has set to $0$. The server
therefore divides it by $1.1$, a heuristic margin, before taking its
sign, for the reason given with Eq.~\eqref{eq:bound}. The sign gives the step vector
\begin{equation}
  \label{eq:step}
  \mathsf{ste} = \tfrac{1}{2}\Bigl( 1 + \mathsf{Sgn}\bigl(
    \tfrac{1}{1.1} (\mathbf{b} - \mathbf{a}) \bigr) \Bigr),
\end{equation}
and the maximum of two label vectors is
\begin{equation}
  \label{eq:max}
  \mathsf{Max}(\mathbf{a}, \mathbf{b}) = \mathbf{a}
    + (\mathbf{b} - \mathbf{a}) \odot \mathsf{ste} .
\end{equation}
In a slot where $\mathbf{b}$ is larger, $\mathsf{ste}$ is $1$ and
$\mathsf{Max}$ returns $\mathbf{b}$. Where $\mathbf{b}$ is smaller,
$\mathsf{ste}$ is $0$ and $\mathsf{Max}$ keeps $\mathbf{a}$. Where the two
are equal, $\mathsf{Max}$ returns their common value whatever
$\mathsf{ste}$ is.

\paragraph{\emph{Core-to-Core Rounds}}
A label may move from point $i+k$ to point $i$ only if the two are
neighbors and both are core points. That condition is itself a product of
masks,
\begin{equation}
  \label{eq:mask}
  \mathbf{m}_k[i] = \mathsf{adj}_k[i] \cdot \mathsf{Core\_mask}[i+k]
    \cdot \mathsf{Core\_mask}[i] ,
\end{equation}
which is $1$ exactly in the slots whose pair may exchange a label. The
server multiplies the rotated labels $\mathsf{Rot}_k(\boldsymbol{\ell})$
by this mask. The product keeps the candidate label in those slots and is
$0$ everywhere else. All labels are non-negative and $\max(x,0) = x$ for
$x \ge 0$, so the zeroed slots leave the maximum unchanged. The mask
$\mathbf{m}_{n-k}$ of the opposite direction has the same form. A round
gives every point the maximum of its label and its candidates for all
offsets $s \in \{k, n-k\}$, $k = 1, \dots, k_{\max}$, and the server runs
$R$ rounds (Section~\ref{sec:rounds}). The masks do not change from round
to round, so the server computes them once.

\paragraph{\emph{Packed Tree Maximum}}
Taking these maxima one offset at a time costs $2k_{\max}$ calls to
$\mathsf{Max}$ per round, each with its own evaluation of $\mathsf{Sgn}$.
The labels, however, occupy only the first $n$ of the $N/2$ slots of a
ciphertext, and a call to $\mathsf{Max}$ costs the same however many
slots carry data. The server therefore places several candidates in the
unused slots and reduces them with a tree of maxima.

Divide the slots into blocks of $n$, block $t$ holding slots
$tn, \dots, (t+1)n - 1$, and let $\mathbf{e}$ be the plaintext vector that
is $1$ in block $0$ and $0$ elsewhere. For $q$ offsets
$s_1, \dots, s_q$ with masks $\mathbf{w}_1, \dots, \mathbf{w}_q$, the
server keeps the labels $\mathbf{a}$ being updated in block $0$ and moves
the candidate of the $t$-th offset, taken from the labels
$\boldsymbol{\lambda}$, into block $t$:
\begin{equation}
  \label{eq:pack}
  \mathbf{P} = \mathbf{e} \odot \mathbf{a} + \sum_{t=1}^{q}
    \mathsf{Rot}\bigl( \mathbf{w}_t \odot
      \mathsf{Rot}_{s_t}(\boldsymbol{\lambda}),\, -tn \bigr).
\end{equation}
Each candidate is $0$ outside block $0$, so the rotation by $-tn$ moves
it into block $t$ and leaves the other blocks untouched. In the
core-to-core rounds, $\mathbf{a} = \boldsymbol{\lambda} = \boldsymbol{\ell}$
and the masks are those of Eq.~\eqref{eq:mask}.

The server then folds the blocks in half. With $\beta = q+1$ occupied
blocks and $h = \lceil \beta/2 \rceil$, it sets
\[
  \mathbf{P} \gets \mathsf{Max}\bigl( \mathbf{P},\,
    \mathsf{Rot}(\mathbf{P}, hn) \bigr),
\]
which gives block $t$ the maximum of itself and block $t+h$, and repeats
with $\beta \gets h$ until $\beta = 1$. After $\lceil \log_2(q+1) \rceil$ calls to
$\mathsf{Max}$, block $0$ holds the maximum of the labels and all $q$
candidates, and $\mathbf{e} \odot \mathbf{P}$ is the updated label vector.
A block that a fold reads beyond the occupied ones holds $0$ or a maximum
of candidates of the same offsets, so it changes nothing.

The fold reads slot $i + hn$ for every occupied slot $i < \beta n$. If this
ran past the end of the ciphertext, the rotation would wrap around and
combine slots that belong to different points. The fold stays inside the
ciphertext when $(\beta + \lceil \beta/2 \rceil)\, n \le N/2$. The server
therefore takes the largest $\beta$ with this property and splits the offsets
$1, \dots, k_{\max}$ into $g = \lceil k_{\max}/m \rceil$ consecutive groups
$G_1, \dots, G_g$ of at most
\[
  m = \min(k_{\max},\, \beta - 1)
\]
offsets each. The groups are processed one after another, each starting
from the labels the previous group produced. A round consists of a forward
pass over $G_1, \dots, G_g$ with the offsets $k$ and a backward pass over
$G_g, \dots, G_1$ with the offsets $n-k$. When $2k = n$, the two offsets
coincide and the backward pass leaves $k$ out.

A round thus costs at most $2g \lceil \log_2(m+1) \rceil$ calls to
$\mathsf{Max}$ instead of $2k_{\max}$. When $k_{\max} < \beta$, all
offsets fit into one group and a round takes only
$2\lceil \log_2(k_{\max}+1) \rceil$ calls. The gain shrinks as $n$ grows:
for $n = 5000$, only $m = 3$ offsets fit, and a round takes about two
thirds of the $2k_{\max}$ calls. The group size depends only on $n$,
$k_{\max}$ and $N$, so the schedule of operations is fixed by public
values. Taking the candidates of a group at once rather than one after
another changes neither the labels the rounds converge to nor the
guarantee that a label advances at least one hop per round
(Section~\ref{sec:rounds}): labels only grow, so within a round every
point receives, from each of its core neighbors, a candidate at least as
large as the label that neighbor had when the round began.

\paragraph{\emph{Core-to-Border Round}}
One further round hands a label to the non-core points. A non-core point
takes the largest label among the core points it neighbors, and keeps $0$
if it neighbors none, in which case it is noise. The mask has the same
form with the roles of the two ends exchanged,
\begin{equation}
  \label{eq:border}
  \mathbf{b}_k[i] = \mathsf{adj}_k[i] \cdot \mathsf{Core\_mask}[i+k]
    \cdot \bigl( 1 - \mathsf{Core\_mask}[i] \bigr),
\end{equation}
and a single round, with a forward and a backward pass as above,
suffices, since a border point passes its label on to no one. Here the
labels $\mathbf{a}$ being updated are a separate vector $\mathbf{z}$ that
starts at $0$, and the candidates are taken from the core labels,
$\boldsymbol{\lambda} = \boldsymbol{\ell}$. A border point adjacent to
core points of two clusters therefore joins the one with the larger
label, the rule discussed in Section~\ref{sec:dbscan}. Core and non-core
labels occupy disjoint slots, so the two are added at the end.
Algorithm~\ref{alg:lp} collects the stage.

\begin{algorithm}[tp]
\caption{\textsf{Label\_Propagation}}
\label{alg:lp}
\begin{algorithmic}[1]
\Require $\mathsf{adj}_1, \dots, \mathsf{adj}_{k_{\max}}$ and
  $\mathsf{adj}_{n-k} = \mathsf{Rot}_{n-k}(\mathsf{adj}_k)$;
  $\mathsf{Core\_mask}$; public $n$, $N$, $k_{\max}$, $R$
\Ensure label ciphertext, slot $i$ carrying the cluster label of point
  $i$, or $0$ if it is noise
\State $\beta \gets$ the largest integer with
  $(\beta + \lceil \beta/2 \rceil)\, n \le N/2$
\State $m \gets \min(k_{\max}, \beta-1)$;\quad
  $g \gets \lceil k_{\max}/m \rceil$
\State $G_j \gets \{(j-1)m + 1, \dots, \min(jm, k_{\max})\}$ for
  $j = 1, \dots, g$
\For{$k = 1$ \textbf{to} $k_{\max}$ \textbf{and} $s \in \{k, n-k\}$}
  \State $\mathbf{m}_s \gets \mathsf{adj}_s \odot
    \mathsf{Rot}_s(\mathsf{Core\_mask}) \odot \mathsf{Core\_mask}$
    \Comment{Eq.~\eqref{eq:mask}}
  \State $\mathbf{b}_s \gets \mathsf{adj}_s \odot
    \mathsf{Rot}_s(\mathsf{Core\_mask}) \odot (1 - \mathsf{Core\_mask})$
    \Comment{Eq.~\eqref{eq:border}}
\EndFor
\State $\boldsymbol{\ell} \gets \mathsf{Core\_mask} \odot \mathbf{u}$
\For{$r = 1$ \textbf{to} $R$}
  \For{$j = 1, \dots, g$} \Comment{forward pass}
    \State $\boldsymbol{\ell} \gets \Call{TreeMax}{\boldsymbol{\ell},
      \boldsymbol{\ell}, \{ (k, \mathbf{m}_k) : k \in G_j \}}$
  \EndFor
  \For{$j = g, \dots, 1$} \Comment{backward pass}
    \State $\boldsymbol{\ell} \gets \Call{TreeMax}{\boldsymbol{\ell},
      \boldsymbol{\ell}, \{ (n-k, \mathbf{m}_{n-k}) : k \in G_j,\,
      2k < n \}}$
  \EndFor
\EndFor
\State $\mathbf{z} \gets 0$
\For{$j = 1, \dots, g$}
  \State $\mathbf{z} \gets \Call{TreeMax}{\mathbf{z}, \boldsymbol{\ell},
    \{ (k, \mathbf{b}_k) : k \in G_j \}}$
\EndFor
\For{$j = g, \dots, 1$}
  \State $\mathbf{z} \gets \Call{TreeMax}{\mathbf{z}, \boldsymbol{\ell},
    \{ (n-k, \mathbf{b}_{n-k}) : k \in G_j,\, 2k < n \}}$
\EndFor
\State \Return $\boldsymbol{\ell} + \mathbf{z}$
\Statex
\Function{TreeMax}{$\mathbf{a}, \boldsymbol{\lambda},
  \{(s_t, \mathbf{w}_t)\}_{t=1}^{q}$}
  \State $\mathbf{P} \gets \mathbf{e} \odot \mathbf{a}$
  \For{$t = 1$ \textbf{to} $q$} \Comment{candidate of $s_t$ into block $t$}
    \State $\mathbf{P} \gets \mathbf{P} + \mathsf{Rot}\bigl(
      \mathbf{w}_t \odot \mathsf{Rot}_{s_t}(\boldsymbol{\lambda}),\,
      -tn \bigr)$
  \EndFor
  \State $\beta \gets q + 1$
  \While{$\beta > 1$} \Comment{fold the blocks in half}
    \State $h \gets \lceil \beta/2 \rceil$
    \State $\mathbf{P} \gets \mathsf{Max}\bigl( \mathbf{P},
      \mathsf{Rot}(\mathbf{P}, hn) \bigr)$
    \State $\beta \gets h$
  \EndWhile
  \State \Return $\mathbf{e} \odot \mathbf{P}$
\EndFunction
\end{algorithmic}
\end{algorithm}